\color{unidarkgreen}{\subsection{Логика деления стороны ВН (ЛД ВН)}}
\color{black}

\begin{enumerate}[label=\arabic{section}.\arabic{subsection}.\arabic*,  align=left, labelsep=4pt, leftmargin=0pt, itemindent=57pt]

\item
Ввод в работу логики деления высокой стороны трансформатора от ТЗНП осуществляется с помощью уставки \textbf{«ЛД ВН от ТЗНП»}. 

\begin{figure}[H]
\centering
\includegraphics[width=0.88\textwidth,height=0.88\textheight,keepaspectratio]{\fbpath/021.1791 ЛД ВН/img/LDVN.pdf}
\captionof{figure}{Функциональная схема ЛД стороны ВН }\label{pict:LD_VN}
\end{figure}

\item
Оперативный ввод/вывод ЛД высокой стороны трансформатора производится с помощью БУ (блока
управления функциями) \textit{«Вывод ЛД ВН»}. 

\item
Первым действием от пуска ТЗНП ВН через выдержку времени \textbf{«Тср ТЗНП откл СВ(ШСВ) ВН»} логика
деления формирует сигналы предварительного деления с действием на отключение СВ (ШСВ) ВН.

\item
Вторым действием от пуска ТЗНП ВН через выдержку времени \textbf{«Тср ТЗНП откл В ВН»} логика деления
формирует сигналы на отключение выключателей высокой стороны трансформатора.

\small
\setlength{\extrarowheight}{0.05cm} % высота ячеек в таблице
\begin{longtable}
{|>{\centering\arraybackslash}m{4.5cm}|>{\centering\arraybackslash}m{4.5cm}|>{\centering\arraybackslash}m{3cm}|>{\centering\arraybackslash}m{3.5cm}|}
\caption{Перечень уставочных параметров ЛД стороны ВН\hfill\vspace{-0.5\baselineskip}}\label{lo_t:tbl1} \\ 
\hhline{|-|-|-|-|}
\rowcolor{gray!30}
Наименование уставки & Диапазон & Шаг &  Значение по умолчанию \\ 
\hhline{|-|-|-|-|}
\endfirsthead

\caption*{\hspace{3pt}\emph{Продолжение таблицы \ref{lo_t:tbl1}\hfill\vspace{-0.5\baselineskip}}} \\ % сделано по ГОСТ 2.105 п.6.8.7
\hhline{|-|-|-|-|}
\rowcolor{gray!30}
Наименование уставки & Диапазон & Шаг &  Значение по умолчанию \\  
\hhline{|-|-|-|-|}
\endhead
\hhline{|-|-|-|-|}

\endfoot
\endlastfoot

%%%%%%%%%%%%содержимое%%%%%%%%%%%%%%

\end{longtable}
\setlength{\extrarowheight}{0.0cm} % возврат параметра высота ячеек в таблице
\normalsize










































\end{enumerate} 